\documentclass[11pt,a4paper]{article}

% ------------------------------------------------------------
% Codificacao, idioma e tipografia -- preserva o padrao dos Volumes I e II
% ------------------------------------------------------------
\usepackage[utf8]{inputenc}
\usepackage[T1]{fontenc}
\usepackage[brazil]{babel}
\usepackage{lmodern}
\usepackage{microtype}
\usepackage{geometry}
\geometry{top=2.2cm,bottom=2.2cm,left=2.25cm,right=2.25cm}
\usepackage{setspace}
\setstretch{1.07}

% ------------------------------------------------------------
% Matematica e tabelas
% ------------------------------------------------------------
\usepackage{amsmath,amssymb,mathtools}
\usepackage{bm}
\usepackage{dsfont}
\usepackage{array,booktabs,longtable,tabularx,multirow,makecell}
\usepackage{enumitem}
\usepackage{siunitx}
\sisetup{output-decimal-marker={,},group-separator={.}}

% ------------------------------------------------------------
% Graficos e caixas -- mesma paleta dos Volumes I e II
% ------------------------------------------------------------
\usepackage{xcolor}
\definecolor{Navy}{HTML}{18324A}
\definecolor{Blue}{HTML}{2D628F}
\definecolor{Teal}{HTML}{2D7A78}
\definecolor{Gold}{HTML}{B8862C}
\definecolor{Red}{HTML}{9B3A36}
\definecolor{LightBlue}{HTML}{EDF4F9}
\definecolor{LightTeal}{HTML}{EEF7F5}
\definecolor{LightGold}{HTML}{FBF6EA}
\definecolor{LightRed}{HTML}{FAEFEE}
\definecolor{LightGray}{HTML}{F4F5F6}
\definecolor{MidGray}{HTML}{7A838A}

\usepackage{tikz}
\usetikzlibrary{arrows.meta,positioning,calc,fit,backgrounds,shapes.geometric,matrix}
\usepackage[most]{tcolorbox}

\newtcolorbox{intuicao}[1][]{
  enhanced,breakable,colback=LightBlue,colframe=Blue,boxrule=0.6pt,
  arc=1.5mm,left=2mm,right=2mm,top=1.5mm,bottom=1.5mm,
  title=Intuição,#1}
\newtcolorbox{formalizacao}[1][]{
  enhanced,breakable,colback=LightTeal,colframe=Teal,boxrule=0.6pt,
  arc=1.5mm,left=2mm,right=2mm,top=1.5mm,bottom=1.5mm,
  title=Formalização,#1}
\newtcolorbox{cuidado}[1][]{
  enhanced,breakable,colback=LightRed,colframe=Red,boxrule=0.6pt,
  arc=1.5mm,left=2mm,right=2mm,top=1.5mm,bottom=1.5mm,
  title=Precisão conceitual,#1}
\newtcolorbox{fonte}[1][]{
  enhanced,breakable,colback=LightGold,colframe=Gold,boxrule=0.6pt,
  arc=1.5mm,left=2mm,right=2mm,top=1.5mm,bottom=1.5mm,
  title=Base normativa,#1}

% ------------------------------------------------------------
% Cabecalho, hiperlinks, secoes
% ------------------------------------------------------------
\usepackage{fancyhdr}
\pagestyle{fancy}
\fancyhf{}
\fancyhead[L]{\small\color{MidGray}Parametrização matemática -- Volume III}
\fancyhead[R]{\small\color{MidGray}\thepage}
\renewcommand{\headrulewidth}{0.3pt}
\setlength{\headheight}{14pt}

\usepackage{titlesec}
\titleformat{\section}{\Large\bfseries\color{Navy}}{\thesection}{0.75em}{}
\titleformat{\subsection}{\large\bfseries\color{Blue}}{\thesubsection}{0.65em}{}
\titleformat{\subsubsection}{\normalsize\bfseries\color{Teal}}{\thesubsubsection}{0.55em}{}
\titlespacing*{\section}{0pt}{1.8em}{0.7em}
\titlespacing*{\subsection}{0pt}{1.35em}{0.5em}

\usepackage{xurl}
\usepackage[hidelinks]{hyperref}
\hypersetup{
  pdftitle={Parametrizacao matematica de objetos contabeis e tributarios - Volume III},
  pdfsubject={Plano de contas, escrituracao, Livro Razao e interface Excel},
  pdfauthor={Caderno tecnico},
  colorlinks=true,
  linkcolor=Navy,
  urlcolor=Blue,
  citecolor=Teal
}
\usepackage[nameinlink,noabbrev]{cleveref}

% ------------------------------------------------------------
% Comandos canonicos herdados dos Volumes I e II
% ------------------------------------------------------------
\newcommand{\R}{\mathbb{R}}
\newcommand{\N}{\mathbb{N}}
\newcommand{\ind}{\mathds{1}}
\newcommand{\BP}{\mathrm{BP}}
\newcommand{\DRE}{\mathrm{DRE}}
\newcommand{\DFC}{\mathrm{DFC}}
\newcommand{\DVA}{\mathrm{DVA}}
\newcommand{\CBS}{\mathrm{CBS}}
\newcommand{\IBS}{\mathrm{IBS}}
\newcommand{\IRPJ}{\mathrm{IRPJ}}
\newcommand{\CSLL}{\mathrm{CSLL}}
\newcommand{\cT}{\mathcal{T}}
\newcommand{\cB}{\mathcal{B}}
\newcommand{\cG}{\mathcal{G}}
\newcommand{\cA}{\mathcal{A}}
\newcommand{\cC}{\mathcal{C}}
\newcommand{\cD}{\mathcal{D}}
\newcommand{\cH}{\mathcal{H}}
\newcommand{\cS}{\mathcal{S}}
\newcommand{\cQ}{\mathcal{Q}}
\newcommand{\cN}{\mathcal{N}}
\newcommand{\periodo}[1]{I_{#1}:=(#1,#1+1]}
\newcommand{\Prov}{\operatorname{Prov}}
\newcommand{\Dep}{\operatorname{Dep}}
\newcommand{\Adm}{\operatorname{Adm}}

% ------------------------------------------------------------
% Objetos novos do Volume III -- escolhidos para evitar colisoes canonicas
% ------------------------------------------------------------
\newcommand{\Plano}{\mathcal{P}}
\newcommand{\cE}{\mathcal{E}}
\newcommand{\Dia}{\mathsf{Dia}}
\newcommand{\Raz}{\mathsf{Raz}}
\newcommand{\Bal}{\mathsf{Bal}}
\newcommand{\Wb}{\mathcal{W}}
\newcommand{\cod}{\operatorname{cod}}
\newcommand{\nome}{\operatorname{nome}}
\newcommand{\nat}{\operatorname{nat}}
\newcommand{\tipo}{\operatorname{tipo}}
\newcommand{\niv}{\operatorname{niv}}
\newcommand{\pai}{\operatorname{pai}}
\newcommand{\map}{\operatorname{map}}
\newcommand{\ordena}{\operatorname{ord}}
\newcommand{\Enc}{\operatorname{Enc}}
\newcommand{\FK}{\operatorname{FK}}
\newcommand{\PK}{\operatorname{PK}}
\newcommand{\Val}{\operatorname{Val}}
\newcommand{\seed}{\operatorname{seed}}

\setlist[itemize]{leftmargin=1.4em,itemsep=0.25em,topsep=0.35em}
\setlist[enumerate]{leftmargin=1.7em,itemsep=0.3em,topsep=0.4em}

\begin{document}
\pagenumbering{gobble}

% ============================================================
% CAPA
% ============================================================
\begin{titlepage}
\centering
\vspace*{1.2cm}
{\Large\color{Blue}\textsc{Caderno de formalização -- Volume III}}\\[0.50cm]
{\huge\bfseries\color{Navy}Plano de contas, escrituração\\[0.15cm]
e interface Excel}\\[0.65cm]
{\Large\color{Teal}Especificação da microestrutura contábil parametrizada\\
para simulação de Diário, Razão e balancete}\\[1.05cm]

\begin{tcolorbox}[width=0.92\textwidth,colback=LightGray,colframe=MidGray,boxrule=0.5pt,arc=2mm]
\centering
\large
\[
\boxed{
\Plano_t
\longrightarrow
u_t
\xrightarrow{\cE_t}
\Lambda_t
\longrightarrow
(\Dia_t,\Raz_t)
\longrightarrow
b_t
\xrightarrow{\cG}
(\BP,\DRE,\DFC,\DVA)
}
\]
\vspace{-1mm}
\small
Estrutura contábil parametrizada $\to$ escrituração verificável $\to$\\
interface física auditável em workbook Excel.
\end{tcolorbox}

\vfill
{\large Versão de referência: 1º de setembro de 2026}\\[0.15cm]
{\small Documento complementar aos Volumes I e II; não constitui parecer contábil ou jurídico.}
\end{titlepage}

\pagenumbering{roman}
\tableofcontents
\clearpage
\pagenumbering{arabic}

% ============================================================
\section{Escopo e posição do Volume III no sistema canônico}
% ============================================================

Os Volumes I e II já estabelecem a constituição abstrata do simulador. O Volume I define o sistema dinâmico contábil-tributário, seus estados, eventos, demonstrações, operadores contábeis e fiscais e a primeira comparação contrafactual. O Volume II adiciona o regime tributário da entidade, seu perfil factual, o vetor mínimo de dados, o seletor de regras efetivas e a camada de proveniência normativa.

O presente Volume III não altera esses objetos. Seu propósito é abrir uma camada que, nos volumes anteriores, permaneceu deliberadamente comprimida: a camada de lançamentos e razão $\Lambda_t$.

\begin{formalizacao}
O ganho diferencial deste volume é resumido por
\[
\boxed{
\Plano_t,\qquad
\Lambda_t,\qquad
\Dia_t,\qquad
\Raz_t,\qquad
b_t,\qquad
\Wb_t.
}
\]
Esses objetos representam, respectivamente, o plano de contas, a família de lançamentos, a visão cronológica da escrituração, a visão por conta, o balancete periódico e a interface física do sistema em Excel.
\end{formalizacao}

\begin{cuidado}
O foco deste volume é a \textbf{microestrutura contábil}. A camada tributária continua dependendo dos objetos dos Volumes I e II. Em particular, um Livro Razão ou um balancete não deve ser tratado, em geral, como estatística suficiente para reconstruir todos os tributos: atributos transacionais podem ter sido perdidos na agregação contábil.
\end{cuidado}

\subsection{Objetivo computacional}

O artefato desejado não é somente uma coleção de números contábeis sintéticos. O objetivo é especificar um pequeno sistema contábil parametrizável, no qual:
\begin{itemize}
    \item Python atua como motor de geração, validação, reconciliação e versionamento;
    \item Excel atua como interface de inspeção, manipulação e auditoria;
    \item o esquema de dados mantém rastreabilidade até operações, lançamentos e regras;
    \item a estrutura gerada é verificável contra restrições contábeis e, posteriormente, calibrável contra dados empíricos.
\end{itemize}

% ============================================================
\section{Recapitulação integral da estrutura canônica dos Volumes I e II}
% ============================================================

Esta seção fixa a notação que o Volume III deve respeitar. Ela funciona como uma constituição do sistema: símbolos herdados não serão reutilizados com significado novo.

\subsection{Tempo, estado e transição}

Mantém-se
\begin{equation}
\boxed{I_t:=(t,t+1],}
\end{equation}
com $x_t\in\mathsf X$ representando o estado econômico-contábil no início de $I_t$. A dinâmica é
\begin{equation}
\boxed{x_{t+1}=F(x_t,u_t;\vartheta_t),}
\label{eq:canonical-dynamics}
\end{equation}
com
\begin{equation}
\boxed{
\vartheta_t=(\theta_t^{\mathrm{acct}},\Theta_t^{\mathrm{tax}}),
\qquad
\Theta_t^{\mathrm{tax}}=(\theta_{j,t})_{j\in\mathsf J_t}.
}
\label{eq:canonical-rules}
\end{equation}

O Balanço Patrimonial permanece uma projeção do estado:
\begin{equation}
\boxed{\BP_t=\pi_{\BP}(x_t).}
\end{equation}

\subsection{Demonstrações e operadores contábeis}

Os objetos de fluxo mantêm as definições do Volume I:
\begin{align}
L_t&=R_t-K_t-E_t+O_t,\\
\DRE_t&=\cG^{\DRE}(x_t,u_t;\theta_t^{\mathrm{acct}}),\\
\DFC_t&=\cG^{\DFC}(x_t,u_t;\theta_t^{\mathrm{acct}}),\\
\DVA_t&=\cG^{\DVA}(x_t,u_t;\theta_t^{\mathrm{acct}}).
\end{align}
Aqui, $L_t$ continua significando \textbf{resultado contábil} e $R_t$ continua significando \textbf{receitas}. Esses símbolos não podem ser reaproveitados neste volume para ``lançamentos'' ou ``Razão''.

A fatoração abstrata do operador contábil é preservada:
\begin{equation}
\boxed{
\cG^S=\cA^S\circ\mathcal K^S\circ\mathcal M^{\mathrm{acct}}\circ\mathcal R^{\mathrm{acct}},
\qquad S\in\{\BP,\DRE,\DFC,\DVA\}.
}
\label{eq:canonical-G}
\end{equation}
Portanto, $\cA^S$ continua sendo operador de agregação e $\mathcal R^{\mathrm{acct}}$ continua sendo operador de reconhecimento.

\subsection{Operações elementares e coleção de eventos}

A operação externa elementar permanece
\begin{equation}
\boxed{
\cT_{k,t}=(d_{k,t},q_{k,t},v_{k,t},s_{k,t},a_{k,t},p_{k,t},\ell_{k,t},\ldots).
}
\label{eq:canonical-event}
\end{equation}
Em particular, $s_{k,t}$ é a data/instante e $p_{k,t}$ já designa condições de pagamento/recebimento. O Volume II reserva $h_{k,t}$ para tratamento tributário especial.

A família de eventos é
\begin{equation}
\boxed{u_t^{\mathrm{tr}}=(\cT_{k,t})_{k=1}^{n_t},}
\end{equation}
e o sistema completo admite
\begin{equation}
\boxed{u_t=u_t^{\mathrm{tr}}\sqcup u_t^{\mathrm{adj}},}
\end{equation}
com $u_t^{\mathrm{adj}}$ contendo ajustes e eventos internos, como depreciação, provisões e apropriações por competência.

\subsection{Camada tributária canônica}

Para cada tributo $j$, preserva-se
\begin{equation}
\boxed{
\theta_{j,t}=
(\theta_{j,t}^{\mathrm{inc}},\theta_{j,t}^{\mathrm{base}},\theta_{j,t}^{\mathrm{aliq}},
\theta_{j,t}^{\mathrm{cred}},\theta_{j,t}^{\mathrm{apur}},\theta_{j,t}^{\mathrm{pag}},\theta_{j,t}^{\mathrm{exc}},\ldots).
}
\end{equation}
Os operadores tributários permanecem
\begin{align}
B_{j,k,t}&=\cB_j(\cT_{k,t};\theta_{j,t}),\\
\tau_{j,k,t}&=\cA_j(\cT_{k,t};\theta_{j,t}),\\
C_{j,k,t}&=\cC_j(\cT_{k,t};\theta_{j,t}),\\
D_{j,k,t}&=\cD_j(\cT_{k,t};\theta_{j,t}).
\end{align}
Logo, $\cA_j$ continua significando operador de alíquota e não será usado para representar plano de contas.

A apuração não cumulativa simplificada permanece
\begin{equation}
\boxed{S_{j,t}^{\mathrm{apur}}=D_{j,t}-C_{j,t}+A_{j,t}^{\mathrm{tax}}.}
\end{equation}
O operador de reforma continua denotado por $\mathfrak R$ e, portanto, essa letra também não será utilizada para o Livro Razão.

\subsection{Objetos introduzidos no Volume II}

A configuração tributária é
\begin{equation}
\boxed{
ho_t=(\rho_t^{\mathrm{ent}},\rho_t^{\mathrm{IR}},\rho_t^{\mathrm{cons}},\rho_t^{\mathrm{esp}}),}
\end{equation}
o perfil factual da entidade é $\eta_t$, e a admissibilidade de um cenário é testada por
\begin{equation}
\boxed{
\chi_t(\rho;\eta_t,\Theta_t^{\mathrm{tax}})\in\{0,1\}.
}
\end{equation}

A regra efetiva é selecionada por
\begin{equation}
\boxed{
\theta_{j,t}^{\mathrm{eff}}=\mathfrak E_{j,t}(\theta_{j,t},\rho_t,\eta_t),
\qquad
\Theta_t^{\mathrm{eff}}=(\theta_{j,t}^{\mathrm{eff}})_{j\in\mathsf J_t}.
}
\label{eq:effective-rules}
\end{equation}

O vetor mínimo de entrada é
\begin{equation}
\boxed{
\zeta_t=(x_t^{\min},u_t^{\min},\rho_t,\eta_t),
}
\label{eq:zeta-III}
\end{equation}
com minimalidade entendida como suficiência funcional para o operador de interesse. Se $H$ é o operador simulado, uma projeção $\Pi$ é suficiente quando existe $\widetilde H$ tal que
\begin{equation}
H(u_t)=\widetilde H(\Pi(u_t)).
\end{equation}

O pacote prático de fontes empresariais permanece
\begin{equation}
\boxed{
Q_t=(\mathrm{ECD}_t,\mathrm{ECF}_t,\mathrm{DocFisc}_t,\mathrm{Cad}_t,\mathrm{Aux}_t),
}
\end{equation}
com normalização
\begin{equation}
\boxed{\widehat\zeta_t=\cN_t(Q_t).}
\end{equation}

Por fim, mantém-se a proveniência normativa
\begin{equation}
\boxed{
\Prov(p)=(\text{fonte},\text{dispositivo},\text{versão},\text{vigência},\text{data de consulta}).
}
\end{equation}

\subsection{Arquitetura canônica consolidada antes do Volume III}

Antes de abrir a escrituração, os dois volumes podem ser resumidos por
\[
\boxed{
Q_t
\xrightarrow{\cN_t}
\widehat\zeta_t
\xrightarrow{\mathfrak E_t}
\Theta_t^{\mathrm{eff}}
\xrightarrow{H}
Y_t,
}
\]
onde o ramo contábil ainda tratava $\Lambda_t$ como camada intermediária abstrata:
\begin{equation}
\boxed{
u_t
\xrightarrow{\mathcal R^{\mathrm{acct}},\mathcal M^{\mathrm{acct}}}
\Lambda_t
\xrightarrow{\mathcal K^S,\cA^S}
S_t.}
\label{eq:lambda-old}
\end{equation}
É precisamente $\Lambda_t$ que será refinada nas próximas seções.

% ============================================================
\section{Ganho informacional do anexo posterior e reconciliação de notação}
% ============================================================

A discussão posterior aos Volumes I e II acrescentou uma camada de modelagem orientada à construção de uma escrituração sintética verossímil. O ganho não está em alterar a estrutura canônica, mas em detalhar o caminho entre eventos e saldos contábeis.

\subsection{Distinção entre plano de contas, saldos e base tributária}

O anexo posterior distinguiu corretamente três níveis conceituais que devem permanecer separados:
\begin{enumerate}
    \item a \textbf{estrutura das contas};
    \item os \textbf{valores/saldos contabilizados};
    \item a \textbf{base tributável} produzida por regras fiscais.
\end{enumerate}

Neste volume, a primeira camada será representada por $\Plano_t$ e a segunda pelo balancete $b_t$. Assim, em vez de introduzir novos símbolos conflitantes, escrevemos
\begin{equation}
\boxed{
(\Plano_t,b_t,u_t,x_t)
\xrightarrow[\rho_t,\eta_t]{\Theta_t^{\mathrm{eff}},\,\cB_j}
B_{j,t}.
}
\label{eq:plan-bal-tax}
\end{equation}

\begin{cuidado}
A expressão \emph{estrutura de contas} não significa ``conjunto de valores sobre os quais incidem os tributos''. O plano de contas informa como a informação contábil é organizada. A base tributária é um objeto derivado da legislação, do regime e dos fatos relevantes.
\end{cuidado}

\subsection{Correções das colisões introduzidas nas respostas posteriores}

A tabela a seguir transforma as ideias do anexo em notação compatível com os volumes canônicos.

\begin{longtable}{p{0.22\textwidth} p{0.28\textwidth} p{0.40\textwidth}}
\toprule
\textbf{Notação provisória do anexo} & \textbf{Notação canônica no Volume III} & \textbf{Razão da correção}\\
\midrule
$\mathcal A$ para estrutura de contas & $\Plano_t$ & $\cA_j$ já é operador de alíquota e $\cA^S$ é agregação.\\
$L_t$ para lançamentos & $\Lambda_t=(\lambda_{\ell,t})$ & $L_t$ já representa resultado contábil na DRE.\\
$R_t$ para Livro Razão & $\Raz_t$ & $R_t$ já representa receitas; $\mathfrak R$ é operador de reforma.\\
$h_\ell$ para cabeçalho & $H^\lambda_{\ell,t}$ & $h_{k,t}$ já representa tratamento tributário especial.\\
$p_{\ell r}$ para partida & $\psi_{\ell r,t}$ & $p_{k,t}$ já representa condições de pagamento/recebimento na operação.\\
$r_t$ para valores de contas & $b_t$ & o Volume II já utiliza $b_t$ como balancete/agregado contábil.\\
\bottomrule
\end{longtable}

\subsection{Novo encadeamento canônico}

Com a reconciliação, o ramo contábil passa a ser
\begin{equation}
\boxed{
(x_t,u_t,\Plano_t)
\xrightarrow{\cE_t(\cdot;\theta_t^{\mathrm{acct}})}
\Lambda_t
\longrightarrow
(\Dia_t,\Raz_t)
\longrightarrow
b_t
\xrightarrow{\cG}
(\BP_{t+1},\DRE_t,\DFC_t,\DVA_t).
}
\label{eq:new-accounting-spine}
\end{equation}

\begin{intuicao}
O ganho informacional pode ser descrito como uma mudança de resolução. Nos volumes anteriores, $\Lambda_t$ era uma ``caixa-preta'' intermediária. Agora explicitaremos o plano de contas que restringe os códigos admissíveis, a estrutura interna de cada lançamento, suas partidas, as duas visões Diário/Razão e a agregação periódica que produz o balancete.
\end{intuicao}

% ============================================================
\section{Evidência estrutural do mundo real: ECD, SPED e normas de escrituração}
% ============================================================

O objetivo de verossimilhança exige uma referência externa para distinguir convenções inventadas de estruturas efetivamente observadas. Para a camada de escrituração, a fonte estrutural oficial mais útil é a Escrituração Contábil Digital (ECD), componente do Sistema Público de Escrituração Digital (SPED).

\begin{fonte}
A Receita Federal descreve a ECD como o módulo do SPED que permite a escrituração eletrônica do Livro Diário, Livro Razão, Balancetes Diários, Balanços e fichas de lançamento. O Manual da ECD -- Leiaute 9, atualizado em janeiro de 2026, contém os registros de plano de contas, centros de custos, saldos periódicos, lançamentos, partidas e demonstrações \cite{ECD_PAGE,ECD2026}.
\end{fonte}

\subsection{Registros oficiais que orientam a spec}

Os registros mais relevantes para o Volume III são:
\begin{longtable}{p{0.13\textwidth} p{0.28\textwidth} p{0.48\textwidth}}
\toprule
\textbf{Registro} & \textbf{Objeto oficial} & \textbf{Uso na especificação}\\
\midrule
I050 & Plano de Contas & estrutura hierárquica, conta sintética/analítica, natureza, código e conta superior;\\
I051 & Plano de Contas Referencial & mapeamento/de-para entre conta analítica da entidade e estrutura referencial;\\
I075 & Histórico Padronizado & cadastro de históricos reutilizáveis;\\
I100 & Centro de Custos & dimensão auxiliar de classificação;\\
I150/I155 & Saldos Periódicos & período, saldo inicial, débitos, créditos, saldo final e indicador D/C;\\
I200 & Lançamento Contábil & cabeçalho do lançamento: identificador, data, valor e tipo;\\
I250 & Partidas do Lançamento & conta analítica, centro de custos, valor, débito/crédito, documento, histórico e participante;\\
I300/I310 & Balancetes Diários & referência adicional para visão de saldos em granularidade diária;\\
J100/J150 & BP/DRE & evidência de que a escrituração alimenta demonstrações estruturadas.\\
\bottomrule
\end{longtable}

\subsection{O que é evidência normativa e o que é escolha de engenharia}

Nem toda coluna necessária ao nosso simulador deve ser confundida com um campo obrigatório da ECD. Adotaremos a seguinte distinção:
\begin{equation}
\boxed{
\text{schema do simulador}
=
\text{núcleo compatível com ECD}
\;\cup\;
\text{campos de engenharia}.
}
\end{equation}

Campos de engenharia podem incluir, por exemplo, identificadores internos de eventos, semente aleatória, versão do gerador, origem sintética/observada, mapeamento para linhas de demonstrações e natureza normal de saldo. Eles devem ser rotulados explicitamente como extensões do modelo.

\subsection{Validade estrutural versus verossimilhança empírica}

\begin{formalizacao}
Definimos duas condições diferentes:
\[
\boxed{
\text{validade estrutural}
\neq
\text{verossimilhança empírica}.
}
\]
A primeira significa satisfazer o schema, as chaves e os invariantes de escrituração. A segunda significa que planos de contas, frequências de lançamentos, valores e combinações observadas parecem plausíveis para entidades reais do arquétipo escolhido.
\end{formalizacao}

A ECD fornece forte evidência para a primeira condição, mas não identifica, por si só, a distribuição típica de contas e lançamentos de uma empresa de determinado setor. Para a segunda condição serão necessários, em etapa posterior, planos de contas, balancetes, Razões ou exportações de ERP reais e anonimizados, além de exemplos profissionais verificáveis.

% ============================================================
\section[Plano de contas como objeto parametrizado]{Plano de contas como objeto parametrizado $\Plano_t$}
% ============================================================

\subsection{Definição abstrata}

Seja $\mathsf C_t$ o conjunto finito de contas vigentes no período:
\begin{equation}
\boxed{\mathsf C_t=\{a_{1,t},\ldots,a_{p_t,t}\}.}
\end{equation}

Em vez de representar cada conta por uma tupla posicional difícil de manter, definimos o plano de contas como uma estrutura relacional:
\begin{equation}
\boxed{
\Plano_t=
(\mathsf C_t,\cod_t,\nome_t,\nat_t,\tipo_t,\niv_t,\pai_t,\map_t).
}
\label{eq:plan-structure}
\end{equation}

As funções têm os seguintes significados:
\begin{itemize}
    \item $\cod_t(a)$: código único da conta;
    \item $\nome_t(a)$: nome da conta;
    \item $\nat_t(a)$: natureza/grupo contábil;
    \item $\tipo_t(a)\in\{S,A\}$: sintética ou analítica;
    \item $\niv_t(a)\in\N$: nível hierárquico;
    \item $\pai_t(a)$: conta sintética imediatamente superior, quando existente;
    \item $\map_t(a)$: mapeamentos auxiliares, por exemplo para demonstrações ou plano referencial.
\end{itemize}

\begin{fonte}
O registro I050 da ECD utiliza, entre outros, os campos DT\_ALT, COD\_NAT, IND\_CTA, NIVEL, COD\_CTA, COD\_CTA\_SUP e CTA. O manual afirma que a ECD é baseada no plano de contas utilizado habitualmente pela pessoa jurídica e distingue contas sintéticas e analíticas \cite{ECD2026}.
\end{fonte}

\subsection{Contas sintéticas e analíticas}

Definimos
\begin{align}
\mathsf C_t^{S}&:=\{a\in\mathsf C_t:\tipo_t(a)=S\},\\
\mathsf C_t^{A}&:=\{a\in\mathsf C_t:\tipo_t(a)=A\}.
\end{align}

Contas sintéticas agregam subcontas; contas analíticas constituem as folhas operacionais da árvore contábil. Para o núcleo compatível com a ECD, as partidas serão associadas a contas analíticas.

\subsection{Plano como árvore ou floresta enraizada}

A relação $\pai_t$ induz uma estrutura hierárquica. Para toda conta não raiz, exigimos
\begin{equation}
\boxed{
\pai_t(a)\in\mathsf C_t^{S},
\qquad
\niv_t(a)=\niv_t(\pai_t(a))+1.
}
\label{eq:hierarchy-invariant}
\end{equation}

Além disso:
\begin{enumerate}
    \item $\cod_t$ é injetiva;
    \item nenhuma conta pode ser sua própria ancestral;
    \item uma conta de nível superior utilizada como pai deve ser sintética;
    \item uma conta analítica utilizada em partida deve existir no plano vigente na data do lançamento;
    \item mudanças de plano devem preservar rastreabilidade entre versões.
\end{enumerate}

O Manual da ECD registra que o plano com contas sintéticas e analíticas deve conter, no mínimo, quatro níveis para fins da formalidade digital tratada pelo CTG 2001 (R3), e fornece como exemplo a cadeia Ativo $\to$ Ativo Circulante $\to$ Disponível $\to$ Caixa \cite{ECD2026,CFC_NORMAS}.

\subsection{Natureza contábil}

A ECD utiliza códigos de natureza para grupos como Ativo, Passivo, Patrimônio Líquido e Contas de Resultado. O simulador preservará esse núcleo e poderá adicionar um atributo de engenharia
\begin{equation}
\delta^{\mathrm{norm}}_t(a)\in\{D,C\},
\end{equation}
que representa a orientação de saldo normal usada pelo motor para apresentação e validação. Esse atributo não deve ser confundido com COD\_NAT da ECD.

\subsection{Mapeamento para demonstrações}

Defina, para cada demonstração $S$, um mapeamento
\begin{equation}
\boxed{
\map_t^S:\mathsf C_t\longrightarrow \mathsf L^S\cup\{\varnothing\},
}
\label{eq:df-map}
\end{equation}
onde $\mathsf L^S$ é o conjunto de linhas/categorias da demonstração. Assim, o plano de contas se torna a ponte entre partidas e agregações do Volume I.

Essa escolha também permite um mapeamento separado para plano referencial, inspirado no registro I051, sem confundir a conta interna da entidade com uma classificação externa.

% ============================================================
\section[Da operação à escrituração: abertura da camada de lançamentos]{Da operação à escrituração: abertura de $\Lambda_t$}
% ============================================================

\subsection{Operador de escrituração}

Introduzimos explicitamente
\begin{equation}
\boxed{
\cE_t:
(x_t,u_t,\Plano_t;\theta_t^{\mathrm{acct}})
\longmapsto
\Lambda_t.
}
\label{eq:bookkeeping-operator}
\end{equation}

A família de lançamentos do período é
\begin{equation}
\boxed{
\Lambda_t=(\lambda_{\ell,t})_{\ell=1}^{M_t}.
}
\end{equation}

Um evento pode gerar zero, um ou vários lançamentos; vários eventos podem também ser consolidados em um lançamento quando a prática e a regra contábil admitirem. Portanto, não imporemos correspondência $1:1$.

\subsection{Relação de rastreabilidade evento--lançamento}

Defina os índices de eventos $\mathsf K_t:=\{1,\ldots,n_t\}$ e de lançamentos $\mathsf L_t:=\{1,\ldots,M_t\}$. A relação
\begin{equation}
\boxed{
\mathcal V_t\subseteq \mathsf K_t\times\mathsf L_t
}
\label{eq:event-entry-link}
\end{equation}
indica quais lançamentos derivam de quais eventos.

\begin{intuicao}
Essa relação evita uma hipótese artificial de cardinalidade. Uma venda pode gerar, por exemplo, um lançamento de receita e outro de baixa do estoque/CMV; um lançamento de folha pode sintetizar muitos eventos operacionais; um ajuste interno pode não possuir documento fiscal de origem.
\end{intuicao}

\subsection{Estrutura de um lançamento}

Cada lançamento é composto por um cabeçalho e suas partidas:
\begin{equation}
\boxed{
\lambda_{\ell,t}
=
\left(H^\lambda_{\ell,t},\Psi_{\ell,t}\right),
\qquad
\Psi_{\ell,t}=(\psi_{\ell r,t})_{r=1}^{m_{\ell,t}}.
}
\label{eq:entry-object}
\end{equation}

Um cabeçalho mínimo compatível com a ideia do I200 é
\begin{equation}
\boxed{
H^\lambda_{\ell,t}
=
(id_{\ell,t},d^\lambda_{\ell,t},V_{\ell,t},\kappa_{\ell,t},d^{\mathrm{ext}}_{\ell,t},\ldots),
}
\end{equation}
onde $id$ é identificador único, $d^\lambda$ a data de lançamento, $V$ o valor do lançamento, $\kappa$ o tipo e $d^{\mathrm{ext}}$ a data do fato para lançamento extemporâneo quando aplicável.

O Manual da ECD distingue no I200 lançamentos normais, de encerramento de contas de resultado e extemporâneos e exige identificador único, data, valor e indicador de tipo \cite{ECD2026}.

\subsection{Partida contábil}

Uma partida será representada por
\begin{equation}
\boxed{
\psi_{\ell r,t}
=
(c_{\ell r,t},cc_{\ell r,t},v^\lambda_{\ell r,t},\delta_{\ell r,t},
q^{\mathrm{doc}}_{\ell r,t},h^{\mathrm{hist}}_{\ell r,t},a^{\mathrm{part}}_{\ell r,t},\ldots),
}
\label{eq:posting}
\end{equation}
com:
\begin{itemize}
    \item $c_{\ell r,t}$: código da conta analítica;
    \item $cc_{\ell r,t}$: centro de custos, se aplicável;
    \item $v^\lambda_{\ell r,t}>0$: valor da partida;
    \item $\delta_{\ell r,t}\in\{D,C\}$: débito ou crédito;
    \item $q^{\mathrm{doc}}_{\ell r,t}$: referência documental;
    \item $h^{\mathrm{hist}}_{\ell r,t}$: histórico;
    \item $a^{\mathrm{part}}_{\ell r,t}$: participante/contraparte quando necessário.
\end{itemize}

Esses campos correspondem ao núcleo observado no I250, que contém conta analítica, centro de custos opcional, valor, indicador D/C, referência documental, histórico e participante opcional \cite{ECD2026}.

\subsection{Invariante fundamental da partida dobrada}

Para todo lançamento $\ell$:
\begin{equation}
\boxed{
\sum_{r:\delta_{\ell r,t}=D}v^\lambda_{\ell r,t}
=
\sum_{r:\delta_{\ell r,t}=C}v^\lambda_{\ell r,t}
=
V_{\ell,t}.
}
\label{eq:double-entry}
\end{equation}

Esse é um invariante estrutural de primeira ordem para o gerador. O Manual da ECD estabelece que o valor do I200 corresponde à soma das partidas com o mesmo indicador D ou C e contém regras de validação das somas de débito e crédito \cite{ECD2026}.

\subsection{Integridade referencial das partidas}

Para toda partida, exigimos
\begin{align}
c_{\ell r,t}&\in \cod_t(\mathsf C_t^{A}),\\
cc_{\ell r,t}&\in\mathsf{CC}_t\cup\{\varnothing\},
\end{align}
onde $\mathsf{CC}_t$ é o cadastro de centros de custo. As chaves estrangeiras devem ser válidas na data do lançamento.

% ============================================================
\section{Diário, Livro Razão e balancete}
% ============================================================

\subsection{Livro Diário como visão cronológica}

O Livro Diário não introduz nova informação econômica: é uma ordenação de $\Lambda_t$ segundo data e identificador. Defina
\begin{equation}
\boxed{
\Dia_t:=\ordena_{(d^\lambda,id)}(\Lambda_t).
}
\label{eq:journal}
\end{equation}

Essa definição é deliberadamente funcional: o objeto primário permanece $\Lambda_t$; o Diário é uma visão ordenada e apresentável dessa família.

\subsection{Livro Razão como projeção por conta}

Para uma conta analítica $a_i$, definimos
\begin{equation}
\boxed{
\Raz_{i,t}
:=
\ordena_{(d^\lambda,id)}
\left\{
(\lambda_{\ell,t},\psi_{\ell r,t}):
 c_{\ell r,t}=\cod_t(a_i)
\right\}.
}
\label{eq:ledger-account}
\end{equation}

O Livro Razão completo é
\begin{equation}
\boxed{
\Raz_t=(\Raz_{i,t})_{a_i\in\mathsf C_t^{A}}.
}
\end{equation}

\subsection{Saldo assinado interno}

Para cálculo computacional, é conveniente utilizar um saldo assinado. Adotaremos, internamente, a convenção ``devedor positivo'':
\begin{equation}
\boxed{
\widetilde b_{i,r}
=
\widetilde b_{i,r-1}
+
D_{i,r}-C_{i,r},
}
\label{eq:signed-balance}
\end{equation}
onde $D_{i,r}$ e $C_{i,r}$ são os débitos e créditos associados ao passo $r$ na conta $i$.

A apresentação contábil pode então ser codificada por
\begin{equation}
\boxed{
\Enc(\widetilde b)=
\begin{cases}
(|\widetilde b|,D),&\widetilde b\ge 0,\\
(|\widetilde b|,C),&\widetilde b<0.
\end{cases}
}
\label{eq:balance-encoding}
\end{equation}

\begin{cuidado}
A convenção de saldo assinado é uma escolha de engenharia interna. Ela não altera a natureza contábil da conta. A interface Excel deve exibir, quando conveniente, valor absoluto e indicador D/C, compatíveis com a forma utilizada nos saldos periódicos da ECD.
\end{cuidado}

\subsection{Balancete periódico}

Seja $J_t$ um conjunto de subperíodos $\tau$ dentro do horizonte. Para cada conta analítica $a_i$ e período $\tau$, definimos
\begin{equation}
\boxed{
 b_{i,\tau}
=
(s^{\mathrm{ini}}_{i,\tau},\delta^{\mathrm{ini}}_{i,\tau},
D_{i,\tau},C_{i,\tau},
 s^{\mathrm{fin}}_{i,\tau},\delta^{\mathrm{fin}}_{i,\tau}).
}
\label{eq:trial-balance-row}
\end{equation}
O balancete é
\begin{equation}
\boxed{b_t=(b_{i,\tau})_{i,\tau}.}
\end{equation}

O formato é inspirado em I150/I155: período, saldo inicial, indicador D/C, total de débitos, total de créditos, saldo final e indicador D/C. O manual prevê periodicidade no máximo mensal e regras de continuidade entre saldos \cite{ECD2026}.

\subsection{Invariantes de saldo}

No espaço assinado,
\begin{equation}
\boxed{
\widetilde b_{i,\tau}^{\mathrm{fin}}
=
\widetilde b_{i,\tau}^{\mathrm{ini}}
+D_{i,\tau}-C_{i,\tau}.
}
\label{eq:trial-balance-invariant}
\end{equation}
Além disso, para períodos consecutivos,
\begin{equation}
\boxed{
\widetilde b_{i,\tau+1}^{\mathrm{ini}}
=
\widetilde b_{i,\tau}^{\mathrm{fin}},
}
\label{eq:continuity}
\end{equation}
salvo mudança de plano formalmente tratada ou outro evento de transição explicitado.

\subsection{Fechamento do ramo contábil}

A arquitetura do Volume I pode agora ser refinada sem alterar o significado de $\cG^S$:
\begin{equation}
\boxed{
(x_t,u_t,\Plano_t)
\xrightarrow{\cE_t}
\Lambda_t
\xrightarrow{\mathrm{views}}
(\Dia_t,\Raz_t,b_t)
\xrightarrow{\cG^S}
S_t,
\qquad S\in\{\BP,\DRE,\DFC,\DVA\}.
}
\label{eq:accounting-final}
\end{equation}

O ganho é interpretativo e computacional: $\cG^S$ deixa de parecer uma agregação direta dos eventos e passa a operar sobre saldos/partidas produzidos por uma escrituração explicitamente validada.

% ============================================================
\section{Workbook Excel como interface física do modelo lógico}
% ============================================================

\subsection{Modelo lógico e interface física}

A distinção central é
\begin{equation}
\boxed{
\text{modelo lógico}
\neq
\text{interface física}.
}
\end{equation}
O modelo lógico é composto por $\Plano_t$, $u_t$, $\Lambda_t$, $\Raz_t$, $b_t$ e pelos demais objetos canônicos. O workbook Excel é uma representação tabular desses objetos para uso humano.

\begin{fonte}
O próprio Manual da ECD 2026 contém uma seção específica sobre abertura de arquivo ECD no Excel e descreve a importação do arquivo delimitado por ``|'' para colunas, com uso de filtros, somas, cálculos e funções. Isso não define nossa arquitetura, mas fornece evidência direta de interoperabilidade prática entre escrituração digital e Excel \cite{ECD2026}.
\end{fonte}

\subsection{Workbook como coleção de tabelas}

Definimos
\begin{equation}
\boxed{
\Wb_t=
(\mathbf T^{\mathrm{cfg}},\mathbf T^{\mathrm{ent}},\mathbf T^{\mathrm{pc}},
\mathbf T^{\mathrm{ev}},\mathbf T^{\mathrm{lct}},\mathbf T^{\mathrm{part}},
\mathbf T^{\mathrm{raz}},\mathbf T^{\mathrm{bal}},\mathbf T^{\mathrm{map}},\ldots).
}
\label{eq:workbook}
\end{equation}
Cada $\mathbf T$ deve ser implementada como uma tabela Excel nomeada, com colunas tipadas, chaves estáveis e validações explícitas.

\subsection{Classes de abas}

Adotaremos três classes funcionais:
\begin{itemize}
    \item \textbf{entrada/configuração}: editáveis pelo usuário;
    \item \textbf{núcleo gerado}: produzido pelo motor e não editado manualmente em fluxo normal;
    \item \textbf{saída/relatório}: derivado, recalculável e auditável.
\end{itemize}

\begin{longtable}{p{0.18\textwidth} p{0.16\textwidth} p{0.18\textwidth} p{0.31\textwidth}}
\toprule
\textbf{Aba} & \textbf{Objeto} & \textbf{Classe} & \textbf{Função}\\
\midrule
README & metadados & entrada & versão da spec, legenda, regras de edição e trilha de origem;\\
CONFIG & $\Omega^{\mathrm{sim}}$ & entrada & horizonte, semente, moeda, escala e parâmetros computacionais;\\
ENTIDADE & $(\eta_t,\rho_t)$ & entrada & atributos do arquétipo e configuração tributária quando aplicável;\\
\makecell[l]{PLANO\_\\ CONTAS} & $\Plano_t$ & entrada/gerada & cadastro hierárquico das contas;\\
\makecell[l]{CENTROS\_\\ CUSTO} & $\mathsf{CC}_t$ & opcional & cadastro auxiliar compatível com I100;\\
\makecell[l]{PARTICI-\\ PANTES} & cadastro & opcional & contrapartes/participantes quando necessários;\\
HISTORICOS & cadastro & opcional & históricos padronizados, inspirado em I075;\\
EVENTOS & $u_t$ & entrada/gerada & eventos econômicos normalizados;\\
\makecell[l]{LANCA-\\ MENTOS} & $H^\lambda_{\ell,t}$ & gerada & cabeçalhos dos lançamentos;\\
PARTIDAS & $\psi_{\ell r,t}$ & gerada & partidas D/C dos lançamentos;\\
\makecell[l]{VINCULO\\ EVENTO\_LCTO} & $\mathcal V_t$ & gerada & relação muitos-para-muitos evento--lançamento;\\
DIARIO & $\Dia_t$ & saída & visão cronológica da escrituração;\\
RAZAO & $\Raz_t$ & saída & visão por conta, filtrável;\\
BALANCETE & $b_t$ & saída & saldos periódicos e movimentos;\\
\makecell[l]{MAPEAMENTO\\ DF} & $\map_t^S$ & entrada & de-para entre contas e linhas das demonstrações;\\
\makecell[l]{BP/DRE/\\ DFC/DVA} & $S_t$ & saída & demonstrações derivadas;\\
VALIDACOES & testes & saída & invariantes, chaves e reconciliações;\\
\makecell[l]{PROVE-\\ NIENCIA} & $\Prov$ & governança & origem de dados, parâmetros e versões;\\
FISCAL\_* & Volume II & reservado & parâmetros, apuração e cenários tributários em etapa posterior.\\
\bottomrule
\end{longtable}

\subsection{Por que uma única aba RAZAO}

No MVP, não criaremos uma aba por conta. Uma tabela única $\mathbf T^{\mathrm{raz}}$, contendo $\cod_t(a)$, nome, data, lançamento, débito, crédito, histórico e saldo, é mais estável e escalável. Filtros, segmentações ou tabelas dinâmicas do Excel podem produzir a visão de cada conta sem multiplicar o número de abas.

\subsection{Esquema relacional e chaves}

As principais restrições de integridade do workbook são:
\begin{align}
\PK(\mathrm{PLANO\_CONTAS})&=\mathrm{COD\_CTA},\\
\PK(\mathrm{LANCAMENTOS})&=\mathrm{NUM\_LCTO},\\
\FK(\mathrm{PARTIDAS.NUM\_LCTO})&\to\mathrm{LANCAMENTOS.NUM\_LCTO},\\
\FK(\mathrm{PARTIDAS.COD\_CTA})&\to\mathrm{PLANO\_CONTAS.COD\_CTA},\\
\FK(\mathrm{PLANO\_CONTAS.COD\_CTA\_SUP})&\to\mathrm{PLANO\_CONTAS.COD\_CTA}.
\end{align}
Se centros de custos e participantes forem usados, suas chaves também devem ser validadas.

\subsection{Fonte de verdade e editabilidade}

A regra operacional recomendada é:
\begin{equation}
\boxed{
\text{editar entradas}
\longrightarrow
\text{regenerar núcleo}
\longrightarrow
\text{recalcular saídas};
}
\end{equation}
e não editar diretamente Razão, balancete ou demonstrações.

Essa disciplina evita que o workbook se torne um arquivo sem proveniência, no qual uma célula manual altera um relatório sem alterar as partidas que deveriam sustentá-lo.

% ============================================================
\section{Contrato mínimo das tabelas do workbook}
% ============================================================

\subsection{PLANO\_CONTAS}

O núcleo recomendado é:
\begin{longtable}{p{0.18\textwidth} p{0.16\textwidth} p{0.15\textwidth} p{0.34\textwidth}}
\toprule
\textbf{Coluna} & \textbf{Tipo} & \textbf{Origem} & \textbf{Regra}\\
\midrule
DT\_ALT & data & ECD/I050 & data de inclusão/alteração;\\
COD\_NAT & texto/código & ECD/I050 & natureza da conta/grupo;\\
IND\_CTA & enum S/A & ECD/I050 & sintética ou analítica;\\
NIVEL & inteiro & ECD/I050 & nível crescente na hierarquia;\\
COD\_CTA & texto & ECD/I050 & chave primária;\\
COD\_CTA\_SUP & texto/nulo & ECD/I050 & conta sintética superior;\\
CTA & texto & ECD/I050 & nome da conta;\\
\makecell[l]{NAT\_SALDO\\ NORMAL} & enum D/C & engenharia & orientação usada pelo motor;\\
COD\_DF & texto/nulo & engenharia & linha da demonstração de destino;\\
ATIVA & booleano & engenharia & vigência operacional no cenário;\\
ORIGEM & enum & governança & observada, sintética, template ou ajustada.\\
\bottomrule
\end{longtable}

\subsection{LANCAMENTOS}

\begin{longtable}{p{0.18\textwidth} p{0.16\textwidth} p{0.15\textwidth} p{0.34\textwidth}}
\toprule
\textbf{Coluna} & \textbf{Tipo} & \textbf{Origem} & \textbf{Regra}\\
\midrule
NUM\_LCTO & texto & I200 & identificador único;\\
DT\_LCTO & data & I200 & data do lançamento;\\
VL\_LCTO & decimal & I200 & valor do lançamento;\\
IND\_LCTO & enum N/E/X & I200 & normal, encerramento ou extemporâneo;\\
\makecell[l]{DT\_LCTO\\ EXT} & data/nulo & I200 & data do fato em lançamento extemporâneo;\\
ID\_GERACAO & texto & engenharia & lote/execução Python que criou o registro;\\
\makecell[l]{VERSAO\\ REGRA} & texto & engenharia & versão da política contábil implementada.\\
\bottomrule
\end{longtable}

\subsection{PARTIDAS}

\begin{longtable}{p{0.18\textwidth} p{0.16\textwidth} p{0.15\textwidth} p{0.34\textwidth}}
\toprule
\textbf{Coluna} & \textbf{Tipo} & \textbf{Origem} & \textbf{Regra}\\
\midrule
ID\_PARTIDA & texto & engenharia & chave interna da linha;\\
NUM\_LCTO & texto & I200/I250 & chave estrangeira do lançamento;\\
COD\_CTA & texto & I250 & conta analítica válida;\\
COD\_CCUS & texto/nulo & I250 & centro de custos;\\
VL\_DC & decimal & I250 & valor positivo da partida;\\
IND\_DC & enum D/C & I250 & natureza da partida;\\
NUM\_ARQ & texto/nulo & I250 & referência documental;\\
\makecell[l]{COD\_HIST\\ PAD} & texto/nulo & I250 & histórico padronizado;\\
HIST & texto & I250 & histórico completo/complementar;\\
COD\_PART & texto/nulo & I250 & participante;\\
\makecell[l]{ID\_\\ ORIGEM} & texto/nulo & engenharia & ponte para evento/documento fonte.\\
\bottomrule
\end{longtable}

\subsection{BALANCETE}

\begin{longtable}{p{0.18\textwidth} p{0.16\textwidth} p{0.15\textwidth} p{0.34\textwidth}}
\toprule
\textbf{Coluna} & \textbf{Tipo} & \textbf{Origem} & \textbf{Regra}\\
\midrule
DT\_INI & data & I150 & início do período;\\
DT\_FIN & data & I150 & fim do período;\\
COD\_CTA & texto & I155 & conta analítica;\\
COD\_CCUS & texto/nulo & I155 & centro de custos;\\
VL\_SLD\_INI & decimal & I155 & saldo inicial absoluto;\\
IND\_DC\_INI & enum D/C & I155 & orientação do saldo inicial;\\
VL\_DEB & decimal & I155 & débitos do período;\\
VL\_CRED & decimal & I155 & créditos do período;\\
VL\_SLD\_FIN & decimal & I155 & saldo final absoluto;\\
IND\_DC\_FIN & enum D/C & I155 & orientação do saldo final.\\
\bottomrule
\end{longtable}

% ============================================================
\section{Responsabilidades de Python e do Excel}
% ============================================================

\subsection{Python como motor}

O motor Python deve concentrar operações difíceis de auditar manualmente quando repetidas em escala:
\begin{enumerate}
    \item gerar ou importar $\Plano_t$;
    \item gerar eventos $u_t$ segundo o arquétipo escolhido;
    \item aplicar $\cE_t$ e materializar $\Lambda_t$;
    \item validar chaves, hierarquia e partida dobrada;
    \item construir Diário, Razão e balancete;
    \item recalcular demonstrações e testes de consistência;
    \item escrever o workbook e, idealmente, relê-lo para validar a materialização;
    \item registrar semente, versão do código e proveniência.
\end{enumerate}

\subsection{Excel como interface}

Excel deve concentrar funções de inspeção e manipulação controlada:
\begin{itemize}
    \item filtros e segmentações;
    \item edição de parâmetros declaradamente editáveis;
    \item fórmulas simples e verificáveis;
    \item pivôs e reconciliações visuais;
    \item leitura por profissionais de contabilidade sem dependência de Python.
\end{itemize}

Para o MVP, recomenda-se o formato $\texttt{.xlsx}$, sem macros, para reduzir dependências e manter portabilidade. Macros ou $\texttt{.xlsm}$ podem ser avaliados apenas quando houver requisito funcional claro.

\subsection{Configuração de simulação}

Para separar o estado econômico de parâmetros puramente computacionais, introduzimos
\begin{equation}
\boxed{
\Omega^{\mathrm{sim}}
=
(I_{\min},I_{\max},\seed,N_{\mathrm{eventos}},\gamma_{\mathrm{escala}},\mathrm{versao},\ldots).
}
\end{equation}
Esse objeto pertence à aba CONFIG, mas não ao estado econômico $x_t$ nem ao perfil factual $\eta_t$.

% ============================================================
\section{Validações e testes de aceitação do Volume III}
% ============================================================

A spec deve declarar invariantes que possam ser executados automaticamente. O resultado da simulação só é admissível se todos os testes estruturais obrigatórios forem satisfeitos.

\subsection{Validações do plano de contas}

\begin{enumerate}
    \item unicidade de COD\_CTA;
    \item inexistência de ciclos na relação de conta superior;
    \item coerência de NIVEL com COD\_CTA\_SUP;
    \item conta superior sintética;
    \item contas utilizadas em partidas são analíticas e existentes;
    \item mapeamentos para demonstrações referenciam categorias válidas.
\end{enumerate}

\subsection{Validações de lançamentos}

Para cada $\lambda_{\ell,t}$:
\begin{align}
\sum_D VL\_DC&=\sum_C VL\_DC,\\
VL\_LCTO&=\sum_D VL\_DC=\sum_C VL\_DC,\\
NUM\_LCTO&\text{ é único.}
\end{align}

\subsection{Validações do Razão e balancete}

\begin{enumerate}
    \item soma dos débitos/créditos por conta no Razão = totais no balancete;
    \item saldo final calculado = saldo final informado;
    \item continuidade entre períodos consecutivos;
    \item soma global de débitos = soma global de créditos no período;
    \item BP derivado satisfaz a identidade patrimonial, dentro da tolerância numérica definida.
\end{enumerate}

\subsection{Validações de proveniência}

Cada lote sintético deve registrar
\begin{equation}
\boxed{
\Prov^{\mathrm{sim}}
=(\seed,\mathrm{versao\_codigo},\mathrm{versao\_spec},\mathrm{template},\mathrm{data\_geracao}).
}
\end{equation}
Para dados observados, a proveniência deve identificar arquivo, período, versão da extração e transformação aplicada.

% ============================================================
\section{Exemplo end-to-end: de eventos ao Razão e às demonstrações}
% ============================================================

Esta seção ilustra a arquitetura sem introduzir tributação. Considere uma entidade criada com seis contas analíticas relevantes:
\begin{center}
\begin{tabular}{lll}
\toprule
Código & Conta & Natureza normal\\
\midrule
1.1.01 & Caixa & D\\
1.1.02 & Clientes & D\\
1.1.03 & Estoques & D\\
2.3.01 & Capital Social & C\\
3.1.01 & Receita de Vendas & C\\
4.1.01 & Custo das Mercadorias Vendidas & D\\
\bottomrule
\end{tabular}
\end{center}

\subsection{Eventos econômicos}

No período, ocorrem quatro eventos:
\begin{enumerate}
    \item aporte de capital de R\$ 100.000 em caixa;
    \item compra à vista de estoque por R\$ 30.000;
    \item venda a prazo por R\$ 50.000, com custo do estoque vendido de R\$ 20.000;
    \item recebimento de R\$ 30.000 do cliente.
\end{enumerate}

O terceiro evento é útil porque demonstra que um único evento econômico pode demandar dois lançamentos contábeis distintos.

\subsection{Lançamentos e partidas}

\begin{longtable}{p{0.10\textwidth} p{0.16\textwidth} p{0.28\textwidth} p{0.10\textwidth} p{0.18\textwidth}}
\toprule
\textbf{Lcto} & \textbf{Conta} & \textbf{Histórico} & \textbf{D/C} & \textbf{Valor (R\$)}\\
\midrule
L001 & Caixa & integralização de capital & D & 100.000\\
L001 & Capital Social & integralização de capital & C & 100.000\\
L002 & Estoques & compra de mercadorias & D & 30.000\\
L002 & Caixa & compra de mercadorias & C & 30.000\\
L003 & Clientes & venda a prazo & D & 50.000\\
L003 & Receita de Vendas & venda a prazo & C & 50.000\\
L004 & CMV & baixa do custo da venda & D & 20.000\\
L004 & Estoques & baixa do custo da venda & C & 20.000\\
L005 & Caixa & recebimento de cliente & D & 30.000\\
L005 & Clientes & recebimento de cliente & C & 30.000\\
\bottomrule
\end{longtable}

Para cada lançamento, os débitos e créditos são iguais. No agregado:
\[
\sum VL_D=230.000=\sum VL_C.
\]

\subsection{Razão por conta}

Os saldos finais em convenção usual são:
\begin{align*}
\mathrm{Caixa}&=100.000\ D,\\
\mathrm{Clientes}&=20.000\ D,\\
\mathrm{Estoques}&=10.000\ D,\\
\mathrm{Capital\ Social}&=100.000\ C,\\
\mathrm{Receita\ de\ Vendas}&=50.000\ C,\\
\mathrm{CMV}&=20.000\ D.
\end{align*}

\subsection{Balancete e demonstrações}

A soma de saldos devedores é
\[
100.000+20.000+10.000+20.000=150.000,
\]
e a soma de saldos credores é
\[
100.000+50.000=150.000.
\]

A DRE simplificada produz
\[
L_t=50.000-20.000=30.000.
\]
Em uma apresentação econômica simplificada, após reconhecer o resultado do período no patrimônio líquido para fechar a posição patrimonial, os ativos somam R\$ 130.000 e o patrimônio líquido econômico também soma R\$ 130.000:
\[
100.000\;\text{Caixa}+20.000\;\text{Clientes}+10.000\;\text{Estoques}
=
100.000\;\text{Capital}+30.000\;\text{Resultado}.
\]

\begin{formalizacao}
O exemplo verifica a cadeia completa:
\[
\boxed{
\Plano_t
\to
u_t
\to
\Lambda_t
\to
\Raz_t
\to
b_t
\to
(\BP_t,\DRE_t).
}
\]
A futura camada tributária não altera esses lançamentos por definição; ela lê os fatos e a escrituração segundo $\Theta_t^{\mathrm{eff}}$, podendo criar lançamentos tributários adicionais quando a política contábil exigir.
\end{formalizacao}

% ============================================================
\section{Conexão com o simulador tributário dos Volumes I e II}
% ============================================================

A microestrutura contábil não substitui o vetor mínimo $\zeta_t$. Ela acrescenta uma fonte estruturada para recuperar parte de $x_t^{\min}$ e de $u_t^{\min}$ e torna auditável a conexão com as demonstrações.

\subsection{Arquitetura conjunta}

Uma arquitetura compatível com os três volumes é
\begin{equation}
\boxed{
Q_t
\xrightarrow{\cN_t}
(\zeta_t,\Plano_t,\Lambda_t,\Raz_t,b_t)
\xrightarrow{\mathfrak E_t}
\Theta_t^{\mathrm{eff}}
\xrightarrow{H}
Y_t.
}
\label{eq:all-volumes}
\end{equation}

No ramo contábil:
\[
(\Plano_t,u_t)\xrightarrow{\cE_t}\Lambda_t\to\Raz_t\to b_t\to S_t.
\]
No ramo tributário:
\[
(x_t^{\min},u_t^{\min},\rho_t,\eta_t;\Theta_t^{\mathrm{eff}})\xrightarrow{H^{\mathrm{tax}}}Y_t^{\mathrm{tax}}.
\]

\subsection{Por que o Razão não elimina a granularidade fiscal}

Se o operador contábil $\mathcal A^{\mathrm{acct}}$ leva eventos a balancete,
\[
\mathcal A^{\mathrm{acct}}:u_t\mapsto b_t,
\]
pode ocorrer
\[
\mathcal A^{\mathrm{acct}}(u_t)=\mathcal A^{\mathrm{acct}}(u'_t)
\quad\text{e}\quad
H^{\mathrm{tax}}(u_t)\ne H^{\mathrm{tax}}(u'_t).
\]
Nesse caso, nem $b_t$ nem uma visão puramente agregada do Razão identificam o resultado tributário. O workbook deve, portanto, preservar a aba EVENTOS e referências documentais mesmo depois que os lançamentos forem gerados.

\subsection{Ponto de integração futuro}

A primeira integração tributária recomendada é adicionar, sem alterar o núcleo contábil:
\begin{itemize}
    \item FISCAL\_PARAM: parâmetros e proveniência de $\Theta_t^{\mathrm{eff}}$;
    \item FISCAL\_OPERACOES: atributos fiscais consumidos por $\cB_j,\cA_j,\cC_j,\cD_j$;
    \item FISCAL\_APURACAO: $S_{j,t}^{\mathrm{apur}}$, $T_{j,t}^{\mathrm{recolher}}$, $C_{j,t}^{\mathrm{saldo}}$;
    \item CENARIOS: combinações admissíveis de $(\rho,\Theta)$ e saídas contrafactuais.
\end{itemize}

% ============================================================
\section{Camada de verossimilhança e calibração empírica}
% ============================================================

Depois de satisfazer a validade estrutural, o projeto precisará calibrar como uma empresa plausível se comporta. Introduzimos, apenas como camada futura,
\begin{equation}
\boxed{
\Gamma=(\Gamma_{\Plano},\Gamma_u,\Gamma_{\cE}),
}
\end{equation}
onde:
\begin{itemize}
    \item $\Gamma_{\Plano}$ controla arquétipos de planos de contas;
    \item $\Gamma_u$ controla frequência, sazonalidade, valores e dependências entre eventos;
    \item $\Gamma_{\cE}$ controla padrões de escrituração e consolidação compatíveis com a política contábil.
\end{itemize}

Esses parâmetros não devem ser escolhidos apenas por intuição. A etapa empírica deve utilizar planos de contas e Razões reais anonimizados, exportações de ERP ou amostras profissionais verificadas. A ECD fornece o envelope estrutural; o corpus empírico fornece a distribuição observada dentro desse envelope.

\begin{cuidado}
Um gerador pode passar todos os testes de partida dobrada, hierarquia e reconciliação e, ainda assim, produzir uma empresa claramente artificial. Validade matemática e plausibilidade empresarial são critérios diferentes e devem ser avaliados separadamente.
\end{cuidado}

% ============================================================
\section{Especificação do MVP do Volume III}
% ============================================================

\subsection{Escopo obrigatório}

O primeiro MVP deve ser capaz de:
\begin{enumerate}
    \item construir um plano de contas hierárquico coerente;
    \item gerar um conjunto finito de eventos de um arquétipo empresarial simples;
    \item converter eventos em lançamentos e partidas por regras declaradas;
    \item garantir partida dobrada e integridade referencial;
    \item derivar Diário, Razão e balancete;
    \item derivar pelo menos BP e DRE;
    \item materializar tudo em $\texttt{.xlsx}$ com tabelas separadas;
    \item produzir uma aba VALIDACOES sem erros obrigatórios;
    \item reproduzir exatamente o mesmo workbook para a mesma semente e versão de configuração, salvo metadados temporais explicitamente excluídos da comparação.
\end{enumerate}

\subsection{Fora do MVP}

Não são requisitos da primeira versão:
\begin{itemize}
    \item reproduzir integralmente o arquivo textual ECD;
    \item cobrir todas as regras do SPED;
    \item implementar todos os CPCs;
    \item simular todos os regimes tributários;
    \item imitar a distribuição real de qualquer setor sem corpus de calibração;
    \item substituir ERP, software contábil ou PGE da ECD.
\end{itemize}

\subsection{Critério de aceitação}

Definimos um vetor de testes
\begin{equation}
\boxed{
\mathcal V^{\mathrm{MVP}}
=(V_{\Plano},V_{\Lambda},V_{\Raz},V_b,V_{\mathrm{DF}},V_{\mathrm{ref}},V_{\mathrm{rep}}).
}
\end{equation}
O cenário é estruturalmente aceito se
\begin{equation}
\boxed{
\prod_{r} \ind\{V_r=\mathrm{ok}\}=1.
}
\end{equation}

% ============================================================
\section{Arquitetura consolidada do Volume III}
% ============================================================

\begin{center}
\resizebox{0.98\textwidth}{!}{%
\begin{tikzpicture}[
  node distance=7mm and 9mm,
  box/.style={draw,rounded corners=2mm,align=center,minimum height=9mm,inner sep=4pt},
  inp/.style={box,fill=LightBlue,draw=Blue,text width=30mm},
  op/.style={box,fill=LightTeal,draw=Teal,text width=34mm},
  mid/.style={box,fill=LightGray,draw=MidGray,text width=34mm},
  gov/.style={box,fill=LightGold,draw=Gold,text width=34mm},
  arr/.style={-{Latex[length=2.2mm]},thick,draw=Navy}
]
\node[inp] (plan) {$\Plano_t$\\plano de contas};
\node[inp,below=of plan] (u) {$u_t$\\eventos};
\node[op,right=of plan,yshift=-8mm] (e) {$\cE_t$\\escrituração};
\node[mid,right=of e] (lam) {$\Lambda_t$\\lançamentos/partidas};
\node[mid,right=of lam,yshift=9mm] (dia) {$\Dia_t$\\visão cronológica};
\node[mid,right=of lam,yshift=-9mm] (raz) {$\Raz_t$\\visão por conta};
\node[mid,right=of raz] (bal) {$b_t$\\balancete};
\node[op,right=of bal] (g) {$\cG^S$\\agregação};
\node[inp,right=of g] (dfs) {$\BP,\DRE,\DFC,\DVA$};
\node[gov,below=18mm of lam] (wb) {$\Wb_t$\\interface Excel};
\node[gov,below=of g] (tax) {$\zeta_t,\rho_t,\eta_t,\Theta_t^{eff}$\\camada tributária};
\draw[arr] (plan) -- (e);
\draw[arr] (u) -- (e);
\draw[arr] (e) -- (lam);
\draw[arr] (lam) -- (dia);
\draw[arr] (lam) -- (raz);
\draw[arr] (raz) -- (bal);
\draw[arr] (bal) -- (g);
\draw[arr] (g) -- (dfs);
\draw[arr] (lam) -- (wb);
\draw[arr] (raz) -- (wb);
\draw[arr] (bal) -- (wb);
\draw[arr] (u.south) |- (tax.west);
\draw[arr] (bal.south) |- (tax.west);
\end{tikzpicture}}
\end{center}

\begin{formalizacao}
A especificação consolidada pode ser lida como quatro contratos:
\begin{enumerate}
    \item \textbf{contrato estrutural}: $\Plano_t$ define quais contas e relações são admissíveis;
    \item \textbf{contrato de escrituração}: $\cE_t$ transforma eventos em $\Lambda_t$ respeitando partida dobrada;
    \item \textbf{contrato de apresentação}: Diário, Razão, balancete e demonstrações são derivados, não editados independentemente;
    \item \textbf{contrato de interface}: $\Wb_t$ materializa os objetos lógicos com chaves, validações e proveniência explícitas.
\end{enumerate}
\end{formalizacao}

% ============================================================
\section{Roteiro de implementação recomendado}
% ============================================================

\subsection{Fase 0 -- schema e testes}

Implementar apenas tipos, colunas, chaves e validadores; popular com um exemplo determinístico pequeno, como o da Seção anterior.

\subsection{Fase 1 -- gerador contábil determinístico}

Construir templates de eventos e regras de escrituração sem aleatoriedade. O objetivo é validar $\cE_t$, $\Lambda_t$, $\Raz_t$ e $b_t$.

\subsection{Fase 2 -- materialização Excel}

Gerar o workbook, aplicar tabelas nomeadas, validações de dados, congelamento de cabeçalhos, filtros e formatação. Reabrir o arquivo via Python e reconciliar valores.

\subsection{Fase 3 -- geração estocástica controlada}

Adicionar $\Omega^{\mathrm{sim}}$ e distribuições paramétricas, mantendo uma semente explícita e testes de reprodutibilidade.

\subsection{Fase 4 -- calibração empírica}

Estimar $\Gamma$ a partir de fontes reais anonimizadas e introduzir critérios quantitativos de verossimilhança.

\subsection{Fase 5 -- integração tributária}

Conectar $\Theta_t^{\mathrm{eff}}$, $\rho_t$, $\eta_t$ e os operadores dos Volumes I/II ao núcleo contábil já validado.

% ============================================================
\appendix
\section{Tabela canônica de símbolos dos três volumes}
% ============================================================

\small
\begin{longtable}{p{0.16\textwidth} p{0.15\textwidth} p{0.60\textwidth}}
\toprule
\textbf{Símbolo} & \textbf{Origem} & \textbf{Significado canônico}\\
\midrule
$I_t$ & I & período $(t,t+1]$;\\
$x_t$ & I & estado econômico-contábil;\\
$u_t$ & I & família de eventos/operações e ajustes;\\
$\cT_{k,t}$ & I & operação/evento elementar;\\
$\vartheta_t$ & I & regras contábeis e tributárias;\\
$\theta_t^{acct}$ & I & regras contábeis;\\
$\Theta_t^{tax}$ & I & conjunto de regras tributárias;\\
$\cG^S$ & I & operador de demonstração $S$;\\
$\Lambda_t$ & I/III & camada de lançamentos, agora explicitada;\\
$L_t$ & I & resultado contábil, não lançamentos;\\
$R_t$ & I & receitas, não Livro Razão;\\
$\cA_j$ & I & operador de alíquota;\\
$\cA^S$ & I & operador de agregação;\\
$\mathfrak R$ & I & operador de reforma;\\
$\rho_t$ & II & configuração tributária da entidade;\\
$\eta_t$ & II & perfil factual da entidade;\\
$\zeta_t$ & II & vetor mínimo de entrada;\\
$\mathfrak E_{j,t}$ & II & seletor de regra efetiva;\\
$\Prov$ & II & proveniência normativa/dados;\\
$Q_t$ & II & pacote de fontes empresariais;\\
$\Plano_t$ & III & plano de contas parametrizado;\\
$\lambda_{\ell,t}$ & III & lançamento contábil;\\
$H^\lambda_{\ell,t}$ & III & cabeçalho do lançamento;\\
$\psi_{\ell r,t}$ & III & partida contábil;\\
$\mathcal V_t$ & III & vínculo evento--lançamento;\\
$\Dia_t$ & III & visão cronológica da escrituração;\\
$\Raz_t$ & III & Livro Razão, visão por conta;\\
$b_t$ & II/III & balancete/agregado periódico de saldos;\\
$\Wb_t$ & III & workbook Excel, interface física;\\
$\Omega^{sim}$ & III & configuração computacional da simulação;\\
$\Gamma$ & III & camada futura de calibração de verossimilhança.\\
\bottomrule
\end{longtable}
\normalsize

% ============================================================
\section{Matriz de ganho informacional e integridade}
% ============================================================

\begin{longtable}{p{0.24\textwidth} p{0.22\textwidth} p{0.39\textwidth}}
\toprule
\textbf{Tema} & \textbf{Situação antes do Volume III} & \textbf{Resultado consolidado}\\
\midrule
Plano de contas & mencionado ou usado como agregado & objeto $\Plano_t$ com hierarquia, atributos e invariantes;\\
$\Lambda_t$ & camada abstrata & família explícita de lançamentos e partidas;\\
Diário/Razão & referências práticas & visões funcionais bem definidas de $\Lambda_t$;\\
Balancete & agregado $b_t$ no Volume II & schema periódico compatível com saldo inicial, D/C, débitos, créditos e saldo final;\\
Rastreabilidade & eventos e documentos separados & relação $\mathcal V_t$ evento--lançamento e chaves do workbook;\\
Partida dobrada & princípio implícito & invariante computável por lançamento;\\
Excel & interface desejada & objeto $\Wb_t$ com abas, chaves, classes de editabilidade e fonte de verdade;\\
Realismo & requisito qualitativo & separação validade estrutural/verossimilhança empírica e futura $\Gamma$;\\
Integração fiscal & Razão potencialmente usado como entrada & preservada a não identificação tributária do Volume II; eventos granulares continuam disponíveis;\\
Notação & colisões nas respostas posteriores & $\Plano_t,\Lambda_t,\Raz_t,\psi$ preservam símbolos canônicos dos Volumes I/II.\\
\bottomrule
\end{longtable}

% ============================================================
\section{Checklist de implementação do workbook}
% ============================================================

\begin{enumerate}
    \item Definir versão da spec e semente da simulação.
    \item Validar $\Plano_t$ antes de gerar qualquer partida.
    \item Gerar EVENTOS com identificadores únicos.
    \item Aplicar $\cE_t$ e criar LANCAMENTOS/PARTIDAS.
    \item Executar teste de partida dobrada para cada NUM\_LCTO.
    \item Validar todas as chaves estrangeiras.
    \item Derivar DIARIO e RAZAO das partidas, sem edição manual.
    \item Derivar BALANCETE e reconciliar movimentos.
    \item Aplicar MAPEAMENTO\_DF e gerar demonstrações.
    \item Executar identidade patrimonial e reconciliações de fluxo.
    \item Registrar PROVENIENCIA e versões.
    \item Salvar $\texttt{.xlsx}$ e reler para teste de integridade física.
\end{enumerate}

% ============================================================
\section{Siglas adicionais}
% ============================================================

\begin{longtable}{p{0.20\textwidth} p{0.70\textwidth}}
\toprule
\textbf{Sigla} & \textbf{Significado}\\
\midrule
ECD & Escrituração Contábil Digital;\\
ECF & Escrituração Contábil Fiscal;\\
SPED & Sistema Público de Escrituração Digital;\\
CFC & Conselho Federal de Contabilidade;\\
ITG & Interpretação Técnica Geral;\\
CTG & Comunicado Técnico Geral;\\
ERP & Enterprise Resource Planning, sistema integrado de gestão empresarial;\\
MVP & Minimum Viable Product, produto mínimo viável.\\
\bottomrule
\end{longtable}

% ============================================================
\section{Referências normativas, contábeis e operacionais}
% ============================================================

\small
\begin{thebibliography}{99}

\bibitem{ECD_PAGE}
RECEITA FEDERAL / SPED. \textbf{Escrituração Contábil Digital -- ECD}. Página institucional. Disponível em: \url{https://www.gov.br/sped/pt-br/assuntos/escrituracoes-digitais/ecd}. Acesso em: 1 set. 2026.

\bibitem{ECD2026}
RECEITA FEDERAL / SPED. \textbf{Manual de Orientação do Leiaute 9 da Escrituração Contábil Digital (ECD)}. Anexo ao ADE Cofis nº 1/2026. Atualização: janeiro de 2026. Disponível em: \url{https://www.gov.br/sped/pt-br/assuntos/escrituracoes-digitais/ecd/manuais-e-documentos-tecnicos/manual_de_orientacao_da_ecd_leiaute_9_janeiro_2026.pdf}. Acesso em: 1 set. 2026.

\bibitem{CFC_NORMAS}
CONSELHO FEDERAL DE CONTABILIDADE. \textbf{Normas Específicas}. Inclui ITG 2000 (R1) -- Escrituração Contábil e CTG 2001 (R3) -- formalidades da escrituração contábil em forma digital para fins do SPED. Disponível em: \url{https://cfc.org.br/tecnica/normas-brasileiras-de-contabilidade/normas-especificas/}. Acesso em: 1 set. 2026.

\bibitem{ECF_PAGE}
RECEITA FEDERAL / SPED. \textbf{Escrituração Contábil Fiscal -- ECF}. Descreve informações fiscais e operacionais para IRPJ/CSLL e recuperação de saldos e contas da ECD. Disponível em: \url{https://www.gov.br/sped/pt-br/assuntos/escrituracoes-digitais/ecf}. Acesso em: 1 set. 2026.

\bibitem{ECF2026}
RECEITA FEDERAL / SPED. \textbf{Manual da ECF -- Leiaute 12}. Atualização de 23 de julho de 2026. Disponível em: \url{https://www.gov.br/sped/pt-br/assuntos/escrituracoes-digitais/ecf/manuais-e-documentos-tecnicos/manual_ecf_leiaute_12_20_05_2026_ac_2025_sit_esp_2026.pdf}. Acesso em: 1 set. 2026.

\bibitem{CPC00}
COMITÊ DE PRONUNCIAMENTOS CONTÁBEIS. \textbf{CPC 00 (R2) -- Estrutura Conceitual para Relatório Financeiro}. Disponível em: \url{https://www.cpc.org.br/Arquivos/Documentos/573_CPC00%28R2%29.pdf}. Acesso em: 1 set. 2026.

\bibitem{CPC03}
COMITÊ DE PRONUNCIAMENTOS CONTÁBEIS. \textbf{CPC 03 (R2) -- Demonstração dos Fluxos de Caixa}. Disponível em: \url{https://www.cpc.org.br/Arquivos/Documentos/183_CPC_03_R2_rev%2021.pdf}. Acesso em: 1 set. 2026.

\bibitem{CPC09}
COMITÊ DE PRONUNCIAMENTOS CONTÁBEIS. \textbf{CPC 09 (R1) -- Demonstração do Valor Adicionado}. Disponível em: \url{https://www.cpc.org.br/Arquivos/Documentos/632_CPC%2009%20R1.pdf}. Acesso em: 1 set. 2026.

\bibitem{CPC26}
COMITÊ DE PRONUNCIAMENTOS CONTÁBEIS. \textbf{CPC 26 (R1) -- Apresentação das Demonstrações Contábeis}. Disponível em: \url{https://www.cpc.org.br/Arquivos/Documentos/312_CPC_26_R1_rev%2023.pdf}. Acesso em: 1 set. 2026.

\end{thebibliography}
\normalsize

\end{document}
